\documentclass[11pt]{article}
\usepackage[english]{babel}
\usepackage[margin=0.82in]{geometry}
\usepackage{amsmath,amssymb,booktabs,tabularx,array}
\usepackage[hidelinks]{hyperref}

\setlength{\parindent}{0pt}
\setlength{\parskip}{0.55\baselineskip}
\setlength{\emergencystretch}{3em}
\newcolumntype{L}[1]{>{\raggedright\arraybackslash}p{#1}}

\title{Technical Review of\\
\emph{Public Emergency Codes: Evidence Audit, Assumption-Conditional Mortality
Scenarios, and Full-Implementation Potential}}
\author{}
\date{August 8, 2026}

\begin{document}
\maketitle

\section*{Overall assessment}

The manuscript is substantially clearer than a conventional ``lives saved''
projection about the distinction between identified evidence and
assumption-conditional planning.  The full manuscript, equations, appendices
embedded in the main text, executable mortality scripts, JSON inputs and
outputs, source transcriptions, tables, captions, and bibliography were
reviewed.  Both mortality simulations were rerun in an isolated copy of the
workspace and reproduced the archived JSON outputs byte-for-byte.  The
manuscript also compiles successfully as a 13-page PDF.

Major revision is still needed.  The current-deployment result is called a
\emph{net} mortality benefit even though its executable model contains only
nonnegative benefit terms.  Its dominant external 988 association is treated
as fixed despite a published confidence interval.  In the mature model, only
five acute-EMS rows are made mutually exclusive; the remaining call, contact,
request, response, and population denominators can overlap each other and the
acute pathways.  The OHCA row is also embedded in a transported-EMS partition
even though its source denominator includes EMS-treated arrests that are not
transported.  These issues do not invalidate the arithmetic that was run, but
they prevent the reported totals from having the incident-level, net-effect
interpretation claimed in several passages.

\section*{Major technical findings}

\subsection*{1. Definite reporting error: the current-deployment result is gross, not net}

\paragraph{Location.}
``Current-deployment attribution scenario,'' the displayed $L_{988}$ and
$L_{\mathrm{OHCA}}$ equations; ``Results that may be quoted,'' which calls the result
``$53$ net statistical lives''; and
\texttt{pec\_evidence\_simulation.py}.

\paragraph{Problem.}
The executable current result is
\[
L_A=L_{988}+L_{L\text{-only}}+L_{A\text{-only}}+L_{\mathrm{joint}},
\]
and every term is nonnegative by construction.  There is no current-deployment
counterpart to the adverse rates used in the mature model for mistranslation,
misrouting, stale information, distraction, false escalation, or added delay.
Consequently, the simulated $53.4$-life median is a gross positive attribution
scenario, not a net difference between beneficial and adverse mortality
effects.

\paragraph{Why it matters.}
``Net'' implies that harms have been subtracted.  The label makes the current
result look directly comparable with the mature net result even though the two
models use different benefit/harm accounting.  The all-positive construction
also explains the mechanical probability of benefit more directly than the
present discussion acknowledges.

\paragraph{Suggested fix.}
Either add explicit current-deployment adverse pathways and report
benefit-minus-harm, or replace ``net statistical lives'' and similar language
with ``gross attributed beneficial lives'' throughout the abstract, results,
quotable text, tables, and conclusion.  If harms cannot be parameterized,
include a separate current adverse-rate sensitivity rather than silently
calling the positive-only result net.

\subsection*{2. Definite uncertainty omission: the load-bearing 988 association is fixed}

\paragraph{Location.}
The displayed $L_{988}=(4{,}372/2.5)gc$ equation, the ``Current-scenario
inputs'' table, the current-attribution interval, and the input field
\texttt{annual\_988\_external\_association} with value $1748.8$.

\paragraph{Problem.}
The cited \texttt{patel2026} study reports $4{,}372$ fewer suicides than
projected with a $95\%$ confidence interval of $3{,}395$--$5{,}349$ over the
2.5-year post-launch period.  The model annualizes only the point estimate and
treats it as fixed.  The sampled causal fraction $c$ addresses attribution and
possible noncausality; it is not a substitute for uncertainty in the estimated
external association.  Likewise, the PEC share $g$ describes a separate
allocation assumption.

\paragraph{Why it matters.}
The 988 component supplies about $38.5$ of the $55.4$ mean current-attribution
lives, so omitting its estimation uncertainty makes the reported
``parameter/scenario interval'' incomplete.  Readers cannot tell that the
interval is conditional on the external association being exactly $1{,}748.8$
per year.

\paragraph{Suggested fix.}
Draw the external association from the study's reported uncertainty or,
preferably, from its model/bootstrap output, and keep that draw distinct from
$g$ and $c$.  If this cannot be done, relabel the interval as conditional on
the fixed $4{,}372/2.5$ association and state the omitted source-estimate
uncertainty beside the headline result.

\subsection*{3. Definite unresolved overlap: the mature total is not based on unique incidents}

\paragraph{Location.}
``Mature full-implementation model,'' the $A_j$, $\lambda_j^{+}$, and
$\lambda_j^{-}$ equations, the ``Complete mature mortality priors'' table,
and the annual-count equations.

\paragraph{Problem.}
The code correctly partitions the five acute-EMS rows, but the total still adds
the 988, 211, 911 diversion/capacity, acute-EMS, and contacts/passive pathways
without a cross-pathway allocation.  These populations are not disjoint.  For
example, the $240$ million 911-call denominator contains the treated-and-
transported 911 EMS responses used by the acute partition, and passive
detection may lead to 911, 988, 211, or EMS activity that is counted again
downstream.

The unit of analysis also changes across rows: 988 contacts, 211 requests,
911 calls, EMS responses, OHCA cases, and U.S. residents are all multiplied by
quantities described as absolute risk changes ``per affected incident.''
Contacts and calls may repeat within one incident or for one person.  No
deduplication or contact-to-unique-incident conversion is specified.

\paragraph{Why it matters.}
Within-pathway benefit/harm/none outcomes can be mutually exclusive while the
same underlying death transition is still credited in more than one pathway.
The direction of the resulting bias is likely upward, but its magnitude cannot
be recovered from the published model.  The predictive count distribution
also treats heterogeneous administrative records as exchangeable incident
trials.

\paragraph{Suggested fix.}
Define one incident-flow or state-transition model that allocates a unique
event first to 988, 211, 911-only, acute EMS, passive-discovery, or explicit
joint categories, then applies one mortality transition.  At minimum, provide
repeat-contact factors and cross-pathway overlap distributions, and report a
sensitivity that attenuates all potentially duplicated pathways together.
Change ``per incident'' to the actual unit wherever a contact-level or
call-level effect is intended.

\subsection*{4. Definite population mismatch: EMS-treated OHCA is placed inside a transported-response partition}

\paragraph{Location.}
The national anchors; the displayed definition of $r^*_{\mathrm{CA}}$; the
``Current-scenario inputs'' table, where $N_{\mathrm{OHCA}}$ is called a
``transported national anchor''; and the acute annual-count multinomial.

\paragraph{Problem.}
$N_{\mathrm{OHCA}}$ is defined from EMS-treated nontraumatic arrests and the
CARES resuscitation-attempt population, which includes arrests pronounced in
the field.  The shared acute denominator $I_A$, by contrast, is explicitly the
number of treated-\emph{and-transported} 911 EMS responses.  Thus
\[
r^*_{\mathrm{CA}}=(N_{\mathrm{OHCA}}/I_A)r_{\mathrm{CA}}
\]
does not convert a subset of $I_A$ into a categorical share; it embeds a count
from a different population in the transported-response multinomial.

This is also required notation/terminology drift: the symbol is first
described as an EMS-treated OHCA count but is later labeled ``transported.''

\paragraph{Why it matters.}
The expected OHCA arithmetic happens to reduce to
$N_{\mathrm{OHCA}}r_{\mathrm{CA}}u_{\mathrm{CA}}$, but the finite-population
interpretation and the claimed common acute incident partition are false for
nontransported arrests.  The resulting negative covariance with transported
acute categories is therefore imposed between populations that are not
actually mutually exclusive parts of one denominator.

\paragraph{Suggested fix.}
Use a common EMS-assessed/treated denominator for all categories, with
transport status modeled explicitly, or place OHCA in a separate stratum and
specify its overlap with transported acute pathways.  Remove ``transported''
from the $N_{\mathrm{OHCA}}$ descriptor unless the source count is replaced by
a transported-only estimate.

\subsection*{5. The printed probability model is incomplete without the code}

\paragraph{Location.}
``Current-scenario inputs,'' ``Mature full-implementation model,'' and the
displayed affected-count and acute-partition equations.

\paragraph{Problem.}
Three implementation choices that materially define the distributions are
absent or ambiguous in the manuscript:
\begin{enumerate}
\item ``beta-PERT with shape 8'' is not defined.  The code uses
\[
\alpha=1+8\frac{m-a}{b-a},\qquad
\beta=1+8\frac{b-m}{b-a},\qquad
X=a+(b-a)\operatorname{Beta}(\alpha,\beta),
\]
but other PERT conventions use different meanings for the shape parameter.
\item The current inputs are sampled independently.  In the mature model,
pathway-specific relevance, success, and ARR draws are independent except for
the five shared multipliers.  Those independence assumptions are not stated.
\item The first mature-model equation defines $A_j=I_jr_ju_j$ for every pathway, while
the implemented acute expectation is $A_j=I_A\tilde q_ju_j$.  The latter is
shown indirectly but never replaces the original definition.
\end{enumerate}

\paragraph{Why it matters.}
The current and mature interval widths depend on the distribution
parameterization and dependence structure.  A reader following only the
printed equations cannot reconstruct the simulation or know which reported
``parameter'' dependencies are assumed absent.

\paragraph{Suggested fix.}
Print the PERT parameterization, state the complete independence/dependence
structure, and define affected opportunities piecewise:
\[
A_j=
\begin{cases}
I_A\tilde q_ju_j,&j\in\mathcal A,\\
I_jr_ju_j,&j\notin\mathcal A.
\end{cases}
\]
Add sensitivity analyses for plausible correlation among pathway relevance,
implementation success, and effect size, especially where difficult cases may
simultaneously have higher relevance and lower success.

\subsection*{6. Likely evidence-conversion issue and descriptor drift in the time slopes}

\paragraph{Location.}
The derivation of $\beta_{\mathrm{CPR}}$, the displayed distribution for
$\beta_{D\to S}$, and the ``Current-scenario inputs'' table.

\paragraph{Problem.}
The prose correctly calls the CPR slope an ``approximate planning transport''
from adjusted delay-category odds ratios and a marginal same-cohort baseline.
The table later calls it an ``exact conversion from U.S. adjusted odds
ratios.''  That is a conflicting descriptor for the same symbol: the logarithm
calculation is exact only after imposing unreported category-time contrasts
and a constant per-minute slope that the study did not fit.

The access slope has a related problem.  The cited adult response-time studies
use heterogeneous populations, endpoints, time categories, and effect
measures, including relative-risk-type statements.  The manuscript does not
show how they produce a constant
$\operatorname{PERT}(0,0.060,0.120)$ \emph{log-odds} coefficient per minute.

\paragraph{Why it matters.}
These slopes feed implementation-facing logistic formulas, so confusing an
odds ratio, risk ratio, categorical comparison, and constant time slope changes
the implied absolute risk gain.  The access-only component is also the largest
OHCA contribution in the current scenario.

\paragraph{Suggested fix.}
Replace ``exact conversion'' with ``author-specified constant-slope
approximation informed by adjusted category associations.''  For
$\beta_{D\to S}$, give the source estimate, endpoint, effect measure, baseline
risk, mathematical conversion, and uncertainty mapping for each bound.  If a
defensible conversion is unavailable, label the full triplet an
author-specified log-odds sensitivity rather than a calibrated registry range.

\subsection*{7. Definite archive defect: the promised reproducibility release is incomplete}

\paragraph{Location.}
``Reproducibility'' and \texttt{source-materials/README.md}.

\paragraph{Problem.}
The manuscript says that
\texttt{reproducibility-manifest.sha256} records hashes of the papers, scripts,
inputs, and outputs, but that file is absent.  The three simulation scripts do
exist locally and reproduce the mortality outputs.  However, the current
\texttt{.gitignore} begins with a wildcard that ignores Python files,
\texttt{scripts/}, and CSV outputs, and those files are not in the Git index.
Thus the active workspace can run the model, but the versioned archive does not
contain the executable package claimed in the paper.

\paragraph{Why it matters.}
The results cannot be independently reproduced from the tracked release, and
the missing manifest makes it impossible to verify that a recipient has the
exact scripts used for the published JSON.  This directly contradicts the
archive-completeness claim.

\paragraph{Suggested fix.}
Revise \texttt{.gitignore}, add the simulation scripts, bundled proofing script,
and required CSV outputs to the release, generate the SHA-256 manifest only
after all revisions, and document a verification command.  Add a clean-clone
test that compiles the paper, reruns both mortality scripts, compares the JSON,
and verifies the manifest.

\section*{Meaningful editorial and notation findings}

\subsection*{Symbols in the joint OHCA equation are used before definition}

\paragraph{Location.}
The joint $L_{\mathrm{OHCA}}$ equation and the preceding definition of $a_L$,
$a_A$, and $a_J$.

\paragraph{Problem and significance.}
$\pi_L$, $u_L$, $\pi_A$, $u_A$, and $t_L$ first appear in the equation block;
their verbal definitions are deferred to the ``Current-scenario inputs'' table.  The
equation is central to the overlap correction, so the delayed definitions make
it difficult to distinguish prevalence, success, and recovered time.

\paragraph{Suggested fix.}
Define each symbol in the prose immediately before the equations, including
its conditioning population and units.  In particular, state that $\pi_L$ is
transported from general medical 911 calls to OHCA and that $t_L$ is recovered
time to layperson CPR initiation.

\subsection*{The count symbol conflicts with receiver compatibility}

\paragraph{Location.}
The shared-multiplier equation uses $C$ for receiver compatibility, while the
annual-count equations use $C_j$ and $C_0$ for acute category counts.

\paragraph{Problem and significance.}
The quantities have different units and roles but appear in adjacent
implementation equations.  This is avoidable notation overload.

\paragraph{Suggested fix.}
Rename the receiver multiplier $R_{\mathrm{comp}}$ or rename category counts
$M_j$ and $M_0$.

\section*{Proofing sweep}

The bundled \texttt{scripts/proofing\_scan.py} was run on the compiled
manuscript.  Its two semicolon-capitalization hits are false positives caused
by the acronym ``PEC'' in the landscape input table.  It found no verified
duplicated punctuation, malformed repeated-preposition frames, product-name
capitalization errors, or inverse-tangent quotient branch ambiguity.  No inverse
trigonometric expression occurs in the manuscript.

The following high-confidence corrections remain:
\begin{itemize}
\item \textbf{The ``Mature modal-input calculations'' table:} the exact modal adverse total is
$106.03371206$, which rounds to $106.0$, not $106.1$.  The current row is also
arithmetically inconsistent because $2{,}763.2-106.1\neq2{,}657.2$.  Replace
$106.1$ with $106.0$.
\item \textbf{The ``Current-scenario inputs'' table:} replace ``exact conversion from
U.S. adjusted odds ratios'' with wording consistent with the manuscript's own
approximation caveat.
\end{itemize}

\section*{Checks completed without an additional concrete error}

The logistic exponents have consistent units because every time slope is per
minute and is multiplied by minutes.  The signs in the time-recovery equations
are internally consistent: recovered time raises modeled survival odds.  The
language/access overlap construction keeps $a_J\leq\min(a_L,a_A)$ and avoids
adding two separate updates for joint cases.  The mature acute annual-count
code now uses a sequential-binomial construction equivalent to the printed
multinomial partition, followed by mutually exclusive benefit/harm/none
counts.  Apart from the $106.1$ table entry, the checked modal totals,
quantiles, threshold probabilities, and sensitivity summaries agree with the
archived executable outputs.

\section*{External-verification status}

Targeted checks support the manuscript's quoted $4{,}372$ 988 association, the
Nguyen adjusted odds ratios of $0.91$ and $0.73$, the video-dispatch trial's
reported 30-day mortality risk difference, the $0.4$-minute escort-associated
vertical-response difference, and the CARES national count and survival
summary.  These checks do not turn any external association into a PEC effect.

The following remain items for external or new empirical verification and
must retain planning-assumption labels:
\begin{itemize}
\item the cross-endpoint factor $\kappa_A$ and the conversion of
scene-to-patient time into mortality;
\item PEC-specific reach, activation, receiver use, overlap, beneficial ARR,
and adverse ARR for every pathway;
\item unique-incident and repeat-contact factors connecting 988 contacts, 211
requests, 911 calls, EMS responses, and passive detection;
\item national OHCA-specific language and material-access prevalence; and
\item the provenance of the unavailable original PEC binary source documents.
\end{itemize}

\section*{Highest-priority fixes}

\begin{enumerate}
\item Make the current scenario genuinely net or relabel it as a gross
positive attribution, and propagate the uncertainty in the external 988
association.
\item Rebuild the mature total around unique incidents or an explicit
cross-pathway state-transition model; do not rely only on the five-row acute
partition.
\item Separate EMS-treated OHCA from the transported-response denominator or
adopt a genuinely common acute population.
\item Complete the printed probability model with the PERT formula,
piecewise affected-count equation, and dependence assumptions; document or
relabel the time-slope conversions.
\item Repair the tracked reproducibility package and manifest, then regenerate
the affected tables and correct the remaining numerical and bibliographic
defects.
\end{enumerate}

After these changes, the manuscript's strong identification caveats would be
matched more closely by its estimands, incident accounting, uncertainty model,
and executable archive.  Until then, the numerical outputs are reproducible as
the stated scripts run, but the current result should be interpreted as a
gross positive scenario and the mature total as an additive planning model
with unresolved cross-pathway overlap.

\end{document}